\begin{lemma}
    If $\best{q}$ is globally asymptotically stable for the deterministic difference inclusion~\eqref{eq: deterministic mces in stochastic section} with any sequences $(\lrate_k)$ and $(\updateset_k)$ that satisfy \autoref{asp: stochastic approximation step size},
    then for any $\epsilon>0$, the solution generated by MC-OPI has a finite number of upcrossings from $\epsilon/2$ to $\epsilon$.
\end{lemma}

\begin{proof}

First, we observe that, if the perturbed inclusion \autoref{eq: perturbed inclusion} is globally asymptotically stable, then for any $\epsilon > 0$, there exists a $\delta > 0$ such that
\begin{align*}
    q_0 \in \best{q} + \delta \ball 
    \quad \implies \quad
    q_\kpz \in \best{q} + \epsilon \ball \quad \forall \ \k>0 .
\end{align*}

Further, since the sequence $(\lrate_k \update_\kpz w_\kpz)$ converges to zero with probability one (\autoref{thm: effects of noise decrease}),
there exists a time step $K$ such that for all $k>K$
\begin{align}
    \max_{s,a} | \lrate_k \update_k w_k(s,a) | < \delta/2 .
\end{align}

Thus, by \autoref{thm: get epsilon close many times}, there exists a time step $K'>K$ when the solution of MC-OPI satisfies
\begin{align}
    q_{K'} \in \best{q} + \frac{\delta}{2} \ball
\end{align}
and if MC-OPI exits this ball, then it is inside the $\delta$-ball because the noise is bounded by $\delta/2$.

When MC-OPI is inside the $\delta$-ball, it behaves like the perturbed inclusion \autoref{eq: perturbed inclusion}. Thus, it does not exit the $\epsilon$-ball around $\best{q}$.

As a result, for any $\epsilon > 0$, there exists a time step $K' \geq 0$ such that $\|\best{q} - q_k\| < \epsilon$ for all $k \geq K'$.
In other words, MC-OPI converges to $\best{q}$ with probability one.


thus if MC-OPI has infinite number of upcrossings from $\epsilon$ to $\epsilon/2$, then deterministic does as well

but this contradicts the fact that deterministic converges


% This proof closely follows the proof of Proposition~4.1 in \citep{Bertsekas1996Neuro-DynamicProgramming}.
Fix an $\epsilon > 0$ 
such that all points $\best{q} + \epsilon \ball \subseteq \region(\best{\pol})$
Therefore, an update is made with $\best{q}$

We observe that for any upcrossing
\begin{align*}
    \dfrac{\epsilon}{2}
    &\leq \| \best{q} - \qval_{K'} \| 
    - \| \best{q} - \qval_{K} \|
    \\
    &= \| (\onefun - \lrate_{K' - 1} \update_{K' - 1}) (\best{q} - \qval_{K' - 1}) - \lrate_{K' - 1} \update_{K' - 1} w_{K' - 1} \| 
    - \| \best{q} - \qval_{K} \|
    \\
    &\leq \| \best{q} - \qval_{K' - 1} \| + \| \lrate_{K' - 1} \update_{K' - 1} w_{K' - 1} \| 
    - \| \best{q} - \qval_{K} \|
    \\
    &\leq \dots
    \\
    &\leq \| \best{q} - \qval_{K} \| + \sum_{k=K}^{K'-1} \| \lrate_{k} \update_{k} w_{k} \| 
    - \| \best{q} - \qval_{K} \|
    \\
    &= \sum_{k=K}^{K'-1} \| \lrate_{k} \update_{k} w_{k} \|
\end{align*}
Thus, if there were an infinite number of upcrossings, then the sequence $(f_k)$ defined in equation \eqref{eq: definition upcrossing sequence} would not converge, which contradicts \autoref{thm: upcrossing noise}.
As a result, the solution generated by MC-OPI has a finite number of upcrossings from $\epsilon/2$ to $\epsilon$.

% Further, suppose that a solution of MC-OPI has infinite upcrossings from $\epsilon/2$ to $\epsilon$, where $\{K_n, \cdots, K'_n\}$ is the $n^\nth$ such upcrossing.


\begin{align}
    \nonumber
    \| \best{q} - \qval_{K_n + 1} \| - \| \best{q} - \qval_{K_n} \|
    &\leq \| \qval_{K_n + 1} - \qval_{K_n} \|
    \\
    \nonumber
    &= \| \lrate_{K_n} \update_{K_n} ( \qval_{\pol_{K_n}} - \qval_{K_n} + w_{K_n} ) \|
    \\
    % \nonumber
    % &\leq \lrate_{K_n} \| \qval_{\pol_{K_n}} - \qval_{K_n} + w_{K_n} \|
    % \\
    \nonumber
    % &\leq \lrate_{K_n} \| q_\opt - \Bar{q} \| + \lrate_{K_n} \| w_{K_n} \|
    &\leq \| \lrate_{K_n} \update_{K_n}( \qval_{\pol_{K_n}} - \qval_{K_n} ) \| + \| \lrate_{K_n} \update_{K_n} w_{K_n} \|
    % \\
    % \label{eq: bound that we can use later}
    % &\leq \lrate_{K_n} \| \qval_{\pol_{K_n}} - \qval_{K_n} \| + \| \lrate_{K_n} \update_{K_n} w_{K_n} \|
\end{align}
As $K_n$ tends to infinity,
\autoref{asp: stochastic conditions on effective learning rate} implies that $\lrate_{K_n} \update_{K_n}$ tends to zero,
\autoref{thm: distance qpol - q is bounded} shows that $\| \qval_{\pol_{K_n}} - \qval_{K_n} \|$ is bounded by $\| \best{q} - \worst{q} \|$,
and \autoref{thm: effects of noise decrease} shows that $\lrate_{K_n} \update_{K_n} w_{K_n}$ converges to zero with probability one.
As a result, there exists an $N$ such that for all $n>N$
\begin{align*}
    \| \lrate_{K_n} \update_{K_n}( \qval_{\pol_{K_n}} - \qval_{K_n} ) \| + \| \lrate_{K_n} \update_{K_n} w_{K_n} \| \leq \epsilon/4 ,
\end{align*}
and consequently
\begin{align}
    \nonumber
    \| \best{q} - \qval_{K_n} \| 
    &\geq
    \| \best{q} - \qval_{K_n + 1} \|  
    - \| \lrate_{K_n} \update_{K_n}( \qval_{\pol_{K_n}} - \qval_{K_n} ) \|
    - \| \lrate_{K_n} \update_{K_n} w_{K_n} \|
    % - \lrate_{K_n} \| \qval_{\pol_{K_n}} - \qval_{K_n} \| 
    % - \| \lrate_{K_n} \update_{K_n} w_{K_n} \|
    \\
    \nonumber
    &\geq \epsilon/2 - \epsilon/4
    \\
    \label{eq: lower bound on beginning of upcrossing}
    &\geq \epsilon/4 .
\end{align}
% Note that the right-hand side converges to zero as $K_n$ tends to infinity because $\lrate_{K_n}$ tends to zero, $\| \qval_{\pol_{K_n}} - \qval_{K_n} \|$ is bounded by $\| \best{q} - \worst{q} \|$
% (\autoref{thm: distance qpol - q is bounded}) and also because of equation \eqref{eq: second useful insight}.
% As a result, since $\| \best{q} - \qval_{K_n + 1} \| \geq \epsilon/2$, equation \eqref{eq: bound that we can use later} implies that for $n$ large enough we have that
% \begin{align}
%     \nonumber
%     \| \best{q} - \qval_{K_n} \| 
%     &\geq
%     \| \best{q} - \qval_{K_n + 1} \|  
%     - \lrate_{K_n} \| \qval_{\pol_{K_n}} - \qval_{K_n} \| 
%     - \| \lrate_{K_n} \update_{K_n} w_{K_n} \|
%     \\
%     \label{eq: lower bound on beginning of upcrossing}
%     &\geq \epsilon/4
% \end{align}

Further, we also observe that
\begin{align*}
    \dfrac{\epsilon}{2}
    &\leq \| \best{q} - \qval_{K'_n} \| - \| \best{q} - \qval_{K_n} \|
    \\
    &\leq \| \qval_{K'_n} - \qval_{K_n} \|
    \\
    &= \bigg\| \sum_{k = K_n}^{K'_n -1} \lrate_k \update_k (\qval_{\pol_k} - q_k + w_k ) \bigg\|
    \\
    &\leq \sum_{k = K_n}^{K'_n -1} \| \lrate_k \update_k ( \qval_{\pol_k} - q_k ) \| + \bigg\| \sum_{k = K_n}^{K'_n -1} \lrate_k \update_k w_k \bigg\| 
    % \\
    % &\leq \sum_{k = K_n}^{K'_n -1} \lrate_k \| \best{q} - \worst{\qval} \| + \bigg\| \sum_{k = K_n}^{K'_n -1} \lrate_k \update_k w_k \bigg\| 
\end{align*}
Lemma~4.1 in \cite{Bertsekas1996Neuro-DynamicProgramming} implies that the second term in the right-hand side converges to zero with probability one.
% Further, Lemma~4.1 in \cite{Bertsekas1996Neuro-DynamicProgramming} implies that 
% the sequence
% \begin{align}
%     u_K = \sum_{k=0}^{K-1} \chi_\kpz \lrate_k \update_\kpz w_\kpz
% \end{align}
% with
% \begin{align}
%     \chi_k
%     \coloneqq
%     \begin{cases}
%         \ 1 &\text{if } \| \qval_\opt - \qval_k \| \leq \epsilon 
%         \\
%         \ 0 &\text{otherwise}
%     \end{cases}
% \end{align}
% converges with probability 1.
% The convergence of $(u_K)$ implies that
% \begin{align}
%     % \label{eq: second useful insight}
%     % &\lim_{n \to \infty} \lrate_{K_n} \update_{K_n} w_{K_n} = 0
%     % \\
%     % \label{eq: first useful insight}
%     \lim_{n \to \infty} \sum_{k = K_n}^{K'_n - 1} \lrate_k \update_k w_k = 0 .
% \end{align}
% Otherwise, the sum of $\chi_k \lrate_k \update_k w_k$ would not converge.
Thus there exists an $N'$ such that for all $n > N'$
\begin{align*}
    
\end{align*}
and consequently
\begin{align*}
    
\end{align*}

Thus, in the limit it holds that
\begin{align}
    \liminf_{n \to \infty} \sum_{k = K_n}^{K'_n -1} \lrate_k \geq \dfrac{\epsilon}{2 \| \best{q} - \worst{q} \|}
\end{align}


Finally, the equations \eqref{eq: def upcrossing intermediate} and \eqref{eq: lower bound on beginning of upcrossing} then show that for large enough $n$, 
\begin{align}
    \| \best{q} - \qval_k \| \geq \epsilon/4 \quad\quad K_n \leq k \leq K_n' - 1
\end{align}
As a result, it holds that
\begin{align}
    \liminf_{n \to \infty} \sum_{k = K_n}^{K'_n -1} \lrate_k \| q_\opt - q_k \|
    \geq 
    \dfrac{\epsilon^2}{8 \| \qval_\opt - \Bar{\qval} \|}
\end{align}
which, after summing over the infinite upcrossings, implies that
\begin{align}
    \sum_{k = 0}^\infty \lrate_\kpz \update_\kpz \| q_\opt - q_k\| 
    % \geq  \sum_{k = 0}^\infty \lrate_\kpz \update_\kpz \| q_\opt - q_k\|^2
    = \infty
\end{align}
\mytodo{make the sum start once $q$ is larger than $\overline{q}$}
However, this contradicts the boundedness of the solution since 
\begin{align}
    \sum_{k = 0}^\infty \lrate_k \update_k \| q_\opt - q_k\|
    % \leq 
    % \sum_{k = 0}^\infty \lrate_\kpz \update_\kpz \| q_\opt - \overline{q}\|^2
    \leq \infty
\end{align}
Thus, there must be a finite number of upcrossing intervals.
\end{proof}

% % Since $\chi_\kpz$ is a function of the random variables in $\hist_\kpz$ and that the summation runs, $u_\kpz$ is as well
% Note that
% % \begin{align}
% %     \E[ \chi_\kpz \lrate_\kpz w_\kpz \mid \hist_\kpz ] 
% %     =
% %     \chi_\kpz \lrate_\kpz \E[ w_\kpz \mid \hist_\kpz ]
% %     = 
% %     0
% % \end{align}
% % and therefore
% \begin{align}
%     \E[u_\kpo \mid \hist_\kpz] 
%     =
%     \E[u_\kpz + \chi_\kpz \lrate_\kpz w_\kpz \mid \hist_\kpz ] 
%     =
%     u_\kpz + \chi_\kpz \lrate_\kpz \E[ w_\kpz \mid \hist_\kpz ] 
%     =
%     u_\kpz
% \end{align}
% Further, if $\chi_\kpz = 0$, then $\E[\|u_\kpo\|^2 \mid \hist_\kpz] = \|u_\kpz\|^2$. If on the other hand $\chi_\kpz = 1$, we have
% \begin{align}
%     \E[\|u_\kpo\|^2 \mid \hist_\kpz]
%     &=
%     \|u_\kpz\|^2 + 2 u_\kpz^\top \lrate_\kpz \E[w_\kpz \mid \hist_\kpz] + \lrate_\kpz^2 \E[\|w_\kpz\|^2 \mid \hist_\kpz]
%     \\
%     &\leq \|u_\kpz\|^2 + \lrate_\kpz^2 (c_1 + c_2 \|\qval_\kpz\|^2)
% \end{align}
% By taking unconditional expectations and summing over $k$, we obtain for all $k$
% \mytodo{define upper bound on stochastic step sizes}
% \begin{align}
%     \E[ \| u_\kpz \|^2] 
%     \leq (c_1 + c_2 \|\qval_\kpz\|^2) \E\bigg[ \sum_{\kappa = 0}^{\infty} \lrate_\kappa^2 \bigg] 
%     \leq (c_1 + c_2 \|\qval_\kpz\|^2) M
% \end{align}
% and since $\|u_k\| \leq 1 + \|u_k\|^2$, we have
% \begin{align}
%     \sup_{k} \E[\|u_\kpz\|] < \infty .
% \end{align}
% The Martingale Convergence Theorem then shows that the sequence $u_\kpz$ converges with probability 1.